\documentclass[11pt]{article}

\usepackage{amssymb,amsmath,amsfonts,amsthm}
\usepackage{lmodern}
\usepackage{xcolor}
\usepackage[colorlinks=true]{hyperref}
\definecolor{cieblue}{RGB}{31, 119, 180}
\hypersetup{
    citecolor=cieblue,
    allcolors=cieblue
}
\usepackage{enumerate}
\newcommand{\warning}[1]{\textcolor{red}{\textbf{Warning:} #1}}

\setlength{\textwidth}{15.5cm}
\setlength{\textheight}{22cm}
\setlength{\oddsidemargin}{0.25cm}
\setlength{\evensidemargin}{0.25cm}
\setlength{\topmargin}{0cm}

\newcommand{\Z}{\mathbb{Z}}
\newcommand{\R}{\mathbb{R}}
\newcommand{\N}{\mathbb{N}}
\newcommand{\Pic}{\operatorname{Pic}^0}

\newtheorem{question}{Question}
\newtheorem{answer}{Discussion}

\title{Research Discussion Notes:\\Picard Groups, Lattices, and Partitions}
\author{Alaittin K\i rt\i \c{s}o\u{g}lu}
\date{April 2026}

\begin{document}
\maketitle
\tableofcontents

%═══════════════════════════════════════════════════════════════════════════════
\section{Picard Group via the Incidence Matrix}
%═══════════════════════════════════════════════════════════════════════════════

\begin{question}
What happens if we define the Picard group using the incidence matrix rather
than the Laplacian?
\end{question}

\begin{answer}
The incidence matrix $B$ is $n \times m$ (vertices $\times$ edges). The
Laplacian is $\Delta = BB^T$.

If we define the ``edge Picard group'' as $\Z^E / \operatorname{im}(B^T B)$:
$B^T B$ is the $m \times m$ edge Laplacian. Its kernel is the cycle space
(dimension $g = m - n + 1$). Its image is in the cut space (dimension $n-1$).

The nonzero eigenvalues of $B^T B$ and $BB^T$ are identical. But $\Z^E$ has
dimension $m$, while $\operatorname{im}(B^T B)$ has rank $n-1$. So the quotient
$\Z^E / \operatorname{im}(B^T B)$ is infinite --- it has a free part of rank
$m - (n-1) = g + 1$ plus a torsion part isomorphic to $\Pic(G)$.

If we restrict to the cut space, we recover $\Pic(G)$ exactly. Other
constructions using $B$ directly give either trivial groups ($\Z^V /
\operatorname{im}(B) \cong \Z$) or infinite groups ($\Z^E /
\operatorname{im}(B^T)$).

The Laplacian is the only construction that yields a finite nontrivial group
from a graph.
\end{answer}


%═══════════════════════════════════════════════════════════════════════════════
\section{The Quotient $\Z^E / (\text{Cut Lattice} + \text{Cycle Lattice})$}
%═══════════════════════════════════════════════════════════════════════════════

\begin{question}
What about $\Z^E / (\operatorname{im}(B^T) + \ker_{\Z}(B))$?
\end{question}

\begin{answer}
Over $\R$: cut space $\oplus$ cycle space $= \R^E$. Every edge vector
decomposes uniquely.

Over $\Z$: $\operatorname{im}(B^T) + \ker_{\Z}(B)$ might not equal $\Z^E$.
However, for graphs with the standard incidence matrix, the fundamental cycles
and cocycles (with respect to any spanning tree) form a \emph{unimodular} basis
of $\Z^E$. This is because the graphic matroid is regular, and regular matroids
are representable by totally unimodular matrices. Therefore:
\[
  \operatorname{im}(B^T) + \ker_{\Z}(B) = \Z^E.
\]
The quotient is trivial. No new group.

This unimodularity is special to graphic matroids. For non-regular matroids,
the quotient could be nontrivial --- this is related to the work of
Baker--Ding--Zhuang on sixth-root-of-unity matroids.
\end{answer}


%═══════════════════════════════════════════════════════════════════════════════
\section{Chip-Firing as Flow}
%═══════════════════════════════════════════════════════════════════════════════

\begin{question}
Chip-firing at vertex $v$ feels like a flow of all chips at $v$ to its
neighbours equally. Is this correct?
\end{question}

\begin{answer}
Yes, that is exactly what it is. Firing vertex $v$:
\begin{itemize}
  \item $v$ loses $\deg(v)$ chips (one per edge),
  \item each neighbour of $v$ gains 1 chip.
\end{itemize}
It is a unit flow pushed out equally along every edge incident to $v$. The
total flow out of $v$ equals $\deg(v)$. Each edge carries exactly 1 unit.

This is why $\Delta_G = BB^T$: the incidence matrix $B$ routes flow along
edges, and $B^T$ aggregates it at vertices. Chip-firing at $v$ is $B$ applied
to the edge flow pattern ``1 unit on every edge of $v$,'' which is
$B^T(v, \cdot)$.

The key limitation: you can only fire all edges of $v$ simultaneously with
equal flow. You cannot choose to send 2 chips along one edge and 0 along
another. This rigidity is what makes the quotient $\Z^V / \operatorname{im}(\Delta)$
finite --- the chip-firing moves are ``too coarse'' to reach every integer
configuration.
\end{answer}


%═══════════════════════════════════════════════════════════════════════════════
\section{Flow Methods and the Picard Group}
%═══════════════════════════════════════════════════════════════════════════════

\begin{question}
Can we implement flow methods with the Picard group?
\end{question}

\begin{answer}
The connection is direct:
\begin{itemize}
  \item Chip-firing at $v$ = push 1 unit of flow out of $v$ along every edge.
  \item Chip-firing at a set $S$ of vertices = push flow out of every vertex
    in $S$ simultaneously.
  \item The characteristic script $\chi(V_i)$ from Xu's thesis = fire all
    vertices in district $V_i$ = push flow out of an entire district across
    its boundary.
\end{itemize}

From Xu's Corollary 4: firing $\chi(V_1)$ on the partition divisor $D^S$
rotates the ordered partition. The flow interpretation: pushing all flow out
of $V_1$ simultaneously redistributes it to neighbouring districts, which
cyclically relabels the partition.

The chip-firing / flow duality gives us:
\begin{itemize}
  \item \textbf{Single vertex fire} = local flow (one node pushes to
    neighbours).
  \item \textbf{District fire $\chi(V_i)$} = global flow (entire district
    pushes across boundary).
  \item \textbf{Difference $D^S - D^{S'}$} between two partition divisors =
    net flow needed to transform one partition into another (Xu's Corollary 3:
    this equals the boundary edges between merged districts).
\end{itemize}

However, chip-firing moves chips (divisor values), not population. Population
is fixed on nodes. The divisor $D^S$ records boundary degrees, not
populations. To connect to population, one needs a divisor encoding population
information, such as $D_{\mathrm{pop}} = p - \bar{p} \cdot \mathbf{1}$.
\end{answer}


%═══════════════════════════════════════════════════════════════════════════════
\section{Changing the Divisor Definition}
%═══════════════════════════════════════════════════════════════════════════════

\begin{question}
If we change the divisor definition according to population flows, how could
the Picard group change?
\end{question}

\begin{answer}
The Picard group does not change. $\Pic(G) = \Z^V / \operatorname{im}(\Delta_G)$
depends only on the graph $G$, not on what the divisor values ``mean.'' The
group structure is fixed by the Laplacian.

What changes is which divisor you are looking at within the fixed group. $D^S$
and $D_{\mathrm{pop}}$ are two different elements of $\operatorname{Div}(G)$.
Chip-firing acts on both via the same group action.

The key insight: if you start at $D_{\mathrm{pop}} = p - \bar{p} \cdot
\mathbf{1}$ and chip-fire, you trace out the coset $[D_{\mathrm{pop}}]$ in
$\Pic(G)$. Every element of this coset is a divisor equivalent to
$D_{\mathrm{pop}}$ --- a different distribution of population deviations
reachable by chip-firing. The coset has exactly $t(G)$ elements.

\textbf{The question becomes:} within the coset $[D_{\mathrm{pop}}]$, which
elements correspond to balanced partitions?

A balanced partition has $|p(\xi_i) - \bar{p}_{\mathrm{district}}|$ small for
each district. In divisor terms: the divisor values, when summed over each
district, should be close to zero. Different elements of $[D_{\mathrm{pop}}]$
redistribute the population deviations differently --- some redistributions
might create balanced partitions, others will not.

So the problem becomes: navigate within a single coset of $\Pic(G)$ to find
divisors whose district sums are balanced. The MCMC walks within this coset
via chip-firing, and the partition emerges from the divisor values.
\end{answer}


%═══════════════════════════════════════════════════════════════════════════════
\section{Why Chip-Firing Finds Spanning Trees}
%═══════════════════════════════════════════════════════════════════════════════

\begin{question}
Why and how is firing a chip useful in finding spanning trees?
\end{question}

\begin{answer}
Start from the sandpile picture. You have a configuration of chips on
vertices. Fix a sink $q$. Keep toppling unstable vertices (those with chips
$\geq$ degree) until everything is stable. The final stable configuration is
unique --- this is the \textbf{abelian property} of the sandpile (Dhar, 1990).

Each recurrent stable configuration corresponds to exactly one spanning tree.
There are $t(G)$ of each. The bijection is concrete: when you run Dhar's
burning algorithm from the sink $q$, the edges that ``carry the fire'' form
the spanning tree.

You are at stable configuration $c$ (corresponding to tree $T$). You add one
chip to some vertex $v$. This might make $v$ unstable. You topple $v$ ---
sending one chip to each neighbour. This might make a neighbour unstable. You
cascade topplings until stable again. You arrive at a new stable configuration
$c'$ (corresponding to tree $T'$).

The map $c \to c'$ is deterministic. By choosing different vertices $v$ to add
chips to, you systematically explore the space of spanning trees.

The exploration has algebraic structure: adding a chip at $v$ and stabilising
corresponds to adding the group element $[\Delta_G(v, \cdot)]$ in $\Pic(G)$.
The group is abelian, so the order does not matter. After firing vertices
$v_1, v_2, \ldots, v_r$, you have added $[\Delta_G(v_1, \cdot)] + \cdots +
[\Delta_G(v_r, \cdot)]$ in the group. The resulting tree depends only on this
sum, not on the order.
\end{answer}


%═══════════════════════════════════════════════════════════════════════════════
\section{Configurations and the Picard Group}
%═══════════════════════════════════════════════════════════════════════════════

\begin{question}
Can I say each configuration is an element in the Picard group? Also, you
said the final stable configuration is unique. What is the name of this
theorem?
\end{question}

\begin{answer}
Almost. Each \emph{recurrent} stable configuration corresponds to a unique
element of $\Pic(G)$. Not every configuration --- only the recurrent ones.
There are $t(G)$ recurrent configurations, $t(G)$ elements of $\Pic(G)$, and
$t(G)$ spanning trees. These three sets are in bijection.

Non-recurrent (transient) configurations do not correspond to group elements.
They appear only during the transient phase and never return.

The confluence of topplings is called the \textbf{abelian property} of the
sandpile model, proved by Dhar (1990). The key statement: if multiple vertices
are unstable, the order in which you topple them does not affect the final
result. This is why the model is called the \emph{abelian} sandpile model.

The abelian property is also why the recurrent configurations form a group ---
the group operation (add two configurations, then stabilise) is well-defined
because stabilisation is unique regardless of toppling order.

Reference: D.~Dhar, ``Self-organized critical state of sandpile automaton
models,'' \emph{Physical Review Letters}, 1990.
\end{answer}


%═══════════════════════════════════════════════════════════════════════════════
\section{Natural Connection Between Sandpile Group and Spanning Trees}
%═══════════════════════════════════════════════════════════════════════════════

\begin{question}
Is there a natural connection between the sandpile group and the set of
spanning trees other than cardinality? Is it just finding a bijection map, or
is there a deeper connection?
\end{question}

\begin{answer}
Deeper. The bijection is not an accident of counting --- it comes from the
structure of the Laplacian.

The Laplacian connects three objects simultaneously:
\begin{enumerate}
  \item \textbf{Spanning trees:} $\det(\Delta^{(q)}) = t(G)$ counts them
    (Kirchhoff).
  \item \textbf{Sandpile group:} $\Z^{n-1} / \operatorname{im}(\Delta^{(q)})$
    defines it.
  \item \textbf{Chip-firing dynamics:} $\Delta_G(v, \cdot)$ is the toppling
    rule.
\end{enumerate}
These are not three separate facts. They are three faces of one object: the
integer lattice $\operatorname{im}(\Delta^{(q)})$ inside $\Z^{n-1}$.

The reduced Laplacian $\Delta^{(q)}$ is an integer matrix. Its columns
generate a sublattice $L$ of $\Z^{n-1}$. The fundamental domain of $L$ (a
parallelotope) has volume $\det(\Delta^{(q)}) = t(G)$. The lattice points of
$\Z^{n-1}$ inside one fundamental domain are the representatives of
$\Z^{n-1}/L = \Pic(G)$. There are exactly $t(G)$ of them.

Each lattice point in the fundamental domain \emph{is} a chip configuration.
The spanning tree bijection shows that this canonical representative encodes a
spanning tree.

\textbf{The bijection exists because spanning trees are bases of the same
lattice that defines the Picard group.} The lattice is
$\operatorname{im}(\Delta^{(q)})$. Its bases are spanning trees. Its cosets
are group elements. The bijection maps each coset to the basis that defines
its canonical representative.
\end{answer}


%═══════════════════════════════════════════════════════════════════════════════
\section{Partitions in the Lattice}
%═══════════════════════════════════════════════════════════════════════════════

\begin{question}
Then we may need to understand what partitions are in the lattice.
\end{question}

\begin{answer}
A spanning tree $T$ is a basis of the lattice
$\operatorname{im}(\Delta^{(q)})$. It has $n-1$ edges, giving $n-1$
independent lattice vectors.

A partition $\xi$ removes $k-1$ edges from $T$, leaving a forest with $n-k$
edges. That forest is a \textbf{partial basis} --- $n-k$ of the $n-1$ lattice
vectors.

So a partition is a \textbf{rank-$(n-k)$ sublattice} of
$\operatorname{im}(\Delta^{(q)})$. Different partitions correspond to
different sublattices of the same rank. Two spanning trees that give the same
partition share the same $n-k$ sublattice vectors (the within-district edges)
and differ only in the remaining $k-1$ vectors (the boundary edges).

In lattice terms:
\begin{itemize}
  \item $\Pic(G) = \Z^{n-1} / L$ where $L = \operatorname{im}(\Delta^{(q)})$,
    a rank-$(n-1)$ lattice.
  \item A spanning tree = a basis of $L$.
  \item A partition = a rank-$(n-k)$ sublattice of $L$.
  \item The fibre = all bases extending this sublattice.
  \item Within-district moves = change the basis without changing the
    sublattice.
  \item Cross-boundary moves = change the sublattice itself.
\end{itemize}

The partition space is analogous to the Grassmannian of rank-$(n-k)$
sublattices of $L$.
\end{answer}


%═══════════════════════════════════════════════════════════════════════════════
\section{The Lattice Explicitly}
%═══════════════════════════════════════════════════════════════════════════════

\begin{question}
Can you help me understand this lattice? Define it explicitly.
\end{question}

\begin{answer}
Fix a graph $G$ with vertices $\{v_1, \ldots, v_n\}$ and a sink $q = v_n$.
The reduced Laplacian $\Delta^{(q)}$ is the $(n-1) \times (n-1)$ matrix
obtained by deleting the last row and column of $\Delta_G$.

The lattice $L$ is:
\[
  L = \{\Delta^{(q)} z : z \in \Z^{n-1}\} \subset \Z^{n-1}.
\]
The columns of $\Delta^{(q)}$ are the generators. Column $i$ is the chip-fire
vector at vertex $v_i$ (restricted to non-sink vertices).

\textbf{Example: the triangle $C_3$.}

Vertices $\{1, 2, 3\}$, sink $q = 3$. Each vertex has degree 2.
\[
  \Delta^{(3)} = \begin{pmatrix} 2 & -1 \\ -1 & 2 \end{pmatrix}.
\]
Column 1 $= (2, -1)$: fire vertex 1. Column 2 $= (-1, 2)$: fire vertex 2.

The lattice $L$ is all integer combinations of $(2, -1)$ and $(-1, 2)$:
\[
  L = \{a(2, -1) + b(-1, 2) : a, b \in \Z\}.
\]
This is a sublattice of $\Z^2$. Its fundamental domain is a parallelogram
with area $= \det(\Delta^{(3)}) = 4 - 1 = 3 = t(C_3)$.

The three cosets of $L$ in $\Z^2$ are $(0,0) + L$, $(1,0) + L$, and
$(0,1) + L$. These are the three elements of $\Pic(C_3) \cong \Z/3$, each
corresponding to one spanning tree.
\end{answer}


%═══════════════════════════════════════════════════════════════════════════════
\section{Kirchhoff's Proof and the Laplacian--Tree Connection}
%═══════════════════════════════════════════════════════════════════════════════

\begin{question}
How is the reduced Laplacian related to spanning trees? How does Kirchhoff
prove it? Are the submatrices with $n-k$ rows related to partitions?
\end{question}

\begin{answer}
\textbf{Kirchhoff's proof:} The Laplacian factors as $\Delta = BB^T$ where
$B$ is the oriented incidence matrix ($n \times m$). The reduced version:
$\Delta^{(q)} = B^{(q)} (B^{(q)})^T$ where $B^{(q)}$ is $(n-1) \times m$.

By Cauchy--Binet:
\[
  \det(\Delta^{(q)}) = \sum_{S \subset E,\, |S|=n-1} \det(B^{(q)}_S)^2
\]
where $B^{(q)}_S$ is the $(n-1) \times (n-1)$ submatrix using columns indexed
by edge set $S$.

The key fact: $\det(B^{(q)}_S) = \pm 1$ if $S$ is a spanning tree, and $0$
otherwise. This is because: if $S$ contains a cycle, the columns are linearly
dependent; if $S$ is a spanning tree, the columns are independent, and the
determinant is $\pm 1$ because $B$ is totally unimodular. Therefore:
\[
  \det(\Delta^{(q)}) = \sum_{T \text{ spanning tree}} (\pm 1)^2 = t(G).
\]

\textbf{Submatrices and partitions:} Yes. The \textbf{Matrix-Forest Theorem}
(Chebotarev and Shamis, 1997) states: for any subset $S \subset V$ with
$|S| = r$, the $r \times r$ minor of $\Delta_G$ indexed by $S$ equals the
number of spanning forests of $G$ where each component tree contains exactly
one vertex from $V \setminus S$ as its root. So:
\begin{itemize}
  \item $(n-1) \times (n-1)$ minor $\to$ spanning trees (1 component,
    1 root).
  \item $(n-2) \times (n-2)$ minor $\to$ spanning forests with 2 components.
  \item $(n-k) \times (n-k)$ minor $\to$ spanning forests with $k$
    components = $k$-partitions (with designated roots).
\end{itemize}

This means the lattice structure at rank $n-k$ directly encodes partitions,
just as the full rank $n-1$ lattice encodes spanning trees. The determinant
of the $(n-k) \times (n-k)$ submatrix counts the partitions (with roots),
just as $\det(\Delta^{(q)})$ counts spanning trees.

\textbf{Key question:} The Picard group lives at rank $n-1$. Is there a
``partition Picard group'' at rank $n-k$?
\end{answer}


%═══════════════════════════════════════════════════════════════════════════════
\section{The Coset Formulation}
%═══════════════════════════════════════════════════════════════════════════════

\begin{answer}[Summary of the coset formulation]
Combining the insights from the previous sections, we arrive at the following
framework:

\textbf{State:} a divisor $D \in [D_{\mathrm{pop}}]$, the coset of
$D_{\mathrm{pop}} = p - \bar{p} \cdot \mathbf{1}$ in $\Pic(G)$.

\textbf{Move:} chip-fire at vertex $v$. $D \to D - \Delta_G(v, \cdot)$. Stays
in the same coset.

\textbf{Partition extraction:} from $D$, define districts by grouping vertices
according to the population deviation pattern encoded in $D$.

\textbf{Reversibility:} chip-firing with symmetric proposal (fire $v$ or
anti-fire $v$ with equal probability) is reversible on the coset. Uniform
stationary distribution on the coset, for free.

\textbf{Irreducibility:} the walk visits every element of the coset
(chip-firing generates $\Pic(G)$, coset is a single orbit). If the partition
extraction rule is surjective, every partition is reachable.

\textbf{Population balance is built in.} $D_{\mathrm{pop}}$ encodes
population. Every element of $[D_{\mathrm{pop}}]$ is a redistribution of the
same total population deviation. The partition extraction looks for clustered
deviations --- districts where the sum is near zero.

\textbf{Open problems:}
\begin{enumerate}
  \item Define a partition extraction rule that is well-defined and surjective.
  \item Prove that the induced chain on partitions is irreducible.
  \item Show polynomial mixing time.
\end{enumerate}
\end{answer}


%═══════════════════════════════════════════════════════════════════════════════
\section{Submatrices of the Laplacian and Partitions}
%═══════════════════════════════════════════════════════════════════════════════

\begin{question}
Does it matter which $(n-k) \times (n-k)$ submatrix we take to generate a
partition? How is the reduced Laplacian related to spanning trees, and are
smaller submatrices related to partitions?
\end{question}

\begin{answer}
Yes. Different $(n-k) \times (n-k)$ submatrices correspond to different
choices of which $k$ vertices to delete --- and the deleted vertices become the
roots of the $k$ components.

The \textbf{Matrix-Forest Theorem} (Chebotarev and Shamis, 1997) states: for
any subset $S \subset V$ with $|S| = k$, the minor of $\Delta_G$ obtained by
deleting the rows and columns indexed by $S$ has determinant equal to the
number of spanning forests of $G$ with $k$ components, where each component is
rooted at one of the vertices in $S$. So:
\begin{itemize}
  \item $(n-1) \times (n-1)$ minor $\to$ spanning trees (1 root).
  \item $(n-2) \times (n-2)$ minor $\to$ 2-component rooted forests.
  \item $(n-k) \times (n-k)$ minor $\to$ $k$-component rooted forests =
    rooted $k$-partitions.
\end{itemize}

Within one minor, forests are grouped by partition. Each partition $\xi$
appears $\prod_i t(G[\xi_i])$ times (different within-district trees, same
components). Two rooted forests from the same minor give the same partition iff
they have the same connected component vertex sets.

The choice of which $k$ vertices to delete determines which vertices are
roots. Different root choices give different counts and different forest sets.
This is an asymmetry with the tree case, where deleting any single vertex gives
the same count $t(G)$.
\end{answer}


%═══════════════════════════════════════════════════════════════════════════════
\section{The Partition Picard Group $\mathrm{Pic}_k(G)$}
%═══════════════════════════════════════════════════════════════════════════════

\begin{question}
What would happen if we consider $\Z^{n-k} / \text{rank-$(n-k)$ sublattice}$?
\end{question}

\begin{answer}
This defines the \textbf{partition Picard group} --- the analogue of $\Pic(G)$
at rank $n-k$ instead of rank $n-1$.

Fix roots $\{q_1, \ldots, q_k\}$. Delete those $k$ rows and columns from
$\Delta_G$ to get $\Delta^{(\{q_1,\ldots,q_k\})}$, an $(n-k) \times (n-k)$
matrix. Define:
\[
  \mathrm{Pic}_k(G) = \Z^{n-k} /
    \operatorname{im}(\Delta^{(\{q_1,\ldots,q_k\})}).
\]
Size $= \det(\Delta^{(\{q_1,\ldots,q_k\})})$ = number of rooted spanning
forests with $k$ components rooted at $q_1, \ldots, q_k$.

Each element of $\mathrm{Pic}_k(G)$ corresponds to one rooted spanning forest.
Chip-firing makes sense: the columns of $\Delta^{(\{q_1,\ldots,q_k\})}$ are
chip-fire vectors at non-root vertices, where chips arriving at any root vertex
disappear (they are deleted from the lattice). This is the
\textbf{multi-sink sandpile}: each district has a designated sink, and chips
drain at roots.
\end{answer}


%═══════════════════════════════════════════════════════════════════════════════
\section{Multi-Sink Bernardi Bijection}
%═══════════════════════════════════════════════════════════════════════════════

\begin{question}
Could we define a bijection similar to Bernardi's bijection for
$\mathrm{Pic}_k(G)$?
\end{question}

\begin{answer}
The ingredients exist but the construction appears to be new.

For Bernardi's bijection at rank $n-1$: fix root $q$, fix ribbon structure.
Tour the graph, record break points, get a divisor. The divisor class bijects
with a spanning tree.

At rank $n-k$ with $k$ roots $\{q_1, \ldots, q_k\}$: start $k$ simultaneous
tours, one from each root $q_i$, each exploring its own district. Each tour
follows the ribbon structure within its territory. When tours meet at boundary
edges, the break points record the boundary structure. The tours' territories
\emph{are} the partition.

Dhar's burning also generalises: burn from $k$ sinks simultaneously. Vertices
catch fire from whichever sink reaches them first. The fire fronts meet at
district boundaries. The burning process defines both the forest (which edges
carry fire) and the partition (which sink claimed each vertex). This is a
\textbf{Voronoi diagram} on the graph with respect to the $k$ root vertices.

The bijection would be: recurrent configurations of the $k$-sink sandpile
$\leftrightarrow$ rooted spanning forests $\leftrightarrow$ elements of
$\mathrm{Pic}_k(G)$.

As far as we are aware, a published multi-sink Bernardi bijection does not
exist. This could be a new contribution.
\end{answer}


%═══════════════════════════════════════════════════════════════════════════════
\section{Relation Between $\Pic(G)$ and $\mathrm{Pic}_k(G)$}
%═══════════════════════════════════════════════════════════════════════════════

\begin{question}
What is the relation between the two Picard groups?
\end{question}

\begin{answer}
There is a natural surjection $\pi : \Pic(G) \to \mathrm{Pic}_k(G)$.

Given a divisor $D \in \Z^{n-1}$ (chips on all vertices except $q$), project
it to $\Z^{n-k}$ by forgetting the coordinates at $q_2, \ldots, q_k$. This
projection respects chip-firing: if $D \sim D'$ in $\Pic(G)$, then
$\pi(D) \sim \pi(D')$ in $\mathrm{Pic}_k(G)$. Therefore:
\[
  \mathrm{Pic}_k(G) \cong \Pic(G) / \ker(\pi).
\]

The kernel $\ker(\pi)$ consists of divisor classes in $\Pic(G)$ that become
trivial in $\mathrm{Pic}_k(G)$ --- classes that can be resolved by chip-firing
when $k-1$ extra sinks are available to absorb chips.

Combinatorially: $\ker(\pi)$ contains the within-district spanning tree
changes. Two spanning trees in the same $\ker(\pi)$-coset map to the same
rooted $k$-forest under projection.

The sizes satisfy:
\[
  |\Pic(G)| = |\ker(\pi)| \cdot |\mathrm{Pic}_k(G)|,
\]
i.e., $t(G) = |\ker(\pi)| \cdot (\text{number of rooted $k$-forests})$.
\end{answer}


%═══════════════════════════════════════════════════════════════════════════════
\section{Properties of $\mathrm{Pic}_k(G)$}
%═══════════════════════════════════════════════════════════════════════════════

\begin{question}
Does every such quotient make a group? How do we know it is abelian?
\end{question}

\begin{answer}
Every quotient $\Z^{n-k} / \operatorname{im}(\Delta^{(\{q_1,\ldots,q_k\})})$
is automatically a finitely generated abelian group. It is abelian because it
is a quotient of $\Z^{n-k}$, which is abelian --- a quotient of an abelian
group is always abelian.

The group is finite (hence useful) when $\Delta^{(\{q_1,\ldots,q_k\})}$ has
full rank $n-k$. For a connected graph with $n > k$, this is guaranteed.

By the Fundamental Theorem of Finitely Generated Abelian Groups, the Smith
normal form of $\Delta^{(\{q_1,\ldots,q_k\})}$ gives the decomposition:
\[
  \mathrm{Pic}_k(G) \cong \Z/d_1 \times \Z/d_2 \times \cdots \times
    \Z/d_{n-k},
\]
where $d_1 \mid d_2 \mid \cdots \mid d_{n-k}$ are the invariant factors.
\end{answer}


%═══════════════════════════════════════════════════════════════════════════════
\section{Cyclic Factors and Spanning Trees}
%═══════════════════════════════════════════════════════════════════════════════

\begin{question}
What does a cyclic group in the decomposition correspond to in the set of
spanning trees?
\end{question}

\begin{answer}
Each cyclic factor $\Z/d_i$ corresponds to an independent ``direction'' through
spanning tree space via chip-firing.

The group $\Z/d_1 \times \Z/d_2$ is like a $d_1 \times d_2$ grid of
positions. Each position is a spanning tree. Moving one step in the first
direction (adding 1 to the first coordinate) changes the tree by a specific
chip-fire combination. Moving $d_1$ steps returns to the same tree. Moving in
the second direction is an independent chip-fire that cycles through $d_2$
trees before returning.

For the triangle $C_3$: $\Pic(C_3) \cong \Z/3$. One cyclic factor, order 3.
Three spanning trees arranged in a cycle. Chip-firing at vertex 1 moves one
step: $T_1 \to T_2 \to T_3 \to T_1$.

For $\mathrm{Pic}_k(G)$: the same interpretation, but in rooted forest space.
Small factors = short loops through few forests. Large factors = long
explorations through many forests.

The \textbf{mixing time} of a random walk on
$\Z/d_1 \times \cdots \times \Z/d_r$ is dominated by the largest factor:
$t_{\mathrm{mix}} = O(d_{\max}^2 \log r)$. So the largest invariant factor of
$\mathrm{Pic}_k(G)$ controls how fast the chain explores all partitions.
\end{answer}


%═══════════════════════════════════════════════════════════════════════════════
\section{Generators and Navigation}
%═══════════════════════════════════════════════════════════════════════════════

\begin{question}
Then finding a tree corresponding to each cyclic group would be enough to
generate the full group? Is it guaranteed that $n-1$ generators visit all
spanning trees?
\end{question}

\begin{answer}
Yes. $\Pic(G)$ has rank $n-1$ (exactly $n-1$ invariant factors). So $n-1$
generators suffice to reach every spanning tree from any starting tree. Every
element of $\Pic(G)$ is a unique integer combination of these $n-1$
generators. For $\mathrm{Pic}_k(G)$: $n-k$ generators suffice to reach every
rooted $k$-forest.

A generator of $\Z/d_i$ is not a single chip-fire at one vertex. It is a
specific \textbf{integer linear combination} of chip-fires at multiple
vertices --- the Smith normal form mixes the original chip-fire vectors to
align with the cyclic factors. The Smith normal form algorithm computes these
generators: when $U \Delta^{(q)} V = \operatorname{diag}(d_1, \ldots, d_{n-1})$,
the columns of $V^{-1}$ are the aligned generators.

\textbf{This gives an algorithm:}
\begin{enumerate}
  \item Compute $\Delta^{(\{q_1,\ldots,q_k\})}$.
  \item Smith normal form $\to$ invariant factors $d_1, \ldots, d_{n-k}$ and
    transformation matrices.
  \item Extract generators $g_1, \ldots, g_{n-k}$.
  \item Starting from any rooted forest, reach any other by: compute the target
    coordinates, apply the corresponding combination of generators.
\end{enumerate}

No MCMC needed for exploration --- you can jump directly to any forest by
applying the right group element. MCMC is only needed for sampling from a
non-uniform distribution (e.g., favouring population balance).

\textbf{The remarkable fact:} the number of generators ($n-k$) is polynomial
in $n$, but the number of objects reached (rooted $k$-forests) is exponential.
The group structure compresses exponential exploration into
polynomial-dimensional navigation.
\end{answer}


%═══════════════════════════════════════════════════════════════════════════════
\section{Risks and Open Questions}
%═══════════════════════════════════════════════════════════════════════════════

The $\mathrm{Pic}_k(G)$ framework is promising but carries several risks:

\begin{enumerate}
  \item \textbf{Root waste:} $\mathrm{Pic}_k(G)$ indexes rooted forests, not
    unrooted partitions. Two rooted forests with the same districts but
    different roots are distinct elements. A chip-fire could change roots
    without changing the partition.

  \item \textbf{Multi-sink Bernardi bijection:} The explicit bijection between
    $\mathrm{Pic}_k(G)$ and rooted $k$-forests via a multi-source touring
    procedure has not been published. It likely works for planar graphs (where
    the ribbon structure controls the tours) but needs proof.

  \item \textbf{Root dependence:} $\mathrm{Pic}_k(G)$ depends on the choice of
    $k$ root vertices. Different choices give different groups (different sizes,
    different invariant factors). There is no canonical root choice for
    redistricting.

  \item \textbf{Kernel characterisation:} The claim that $\ker(\pi)$
    corresponds exactly to within-district moves needs careful verification.
    Some cross-boundary chip-fires might cancel at the extra sinks.
\end{enumerate}

These are concrete, testable questions that can be resolved by computation on
small examples (triangle, $6 \times 6$ grid) before attempting general proofs.


%═══════════════════════════════════════════════════════════════════════════════
\section{The Structure of $k$-Partitions Inside $k$-Forests}
%═══════════════════════════════════════════════════════════════════════════════

\begin{question}
Can we identify the structure of $k$-partitions inside the space of
$k$-forests, and embed both inside $\Pic(G)$?
\end{question}

\begin{answer}
The group $\mathrm{Pic}_k(G)$ acts on rooted $k$-forests. The rooted
$k$-forests partition into equivalence classes, where each class consists of
all forests giving the same (unrooted) $k$-partition.

\paragraph{The fibre of a partition.}
For a fixed connected $k$-partition $\xi = (\xi_1, \ldots, \xi_k)$, the
number of rooted $k$-forests producing $\xi$ is:
\[
  |\text{fibre}(\xi)| = \prod_{i=1}^{k} |V(\xi_i)| \cdot t(G[\xi_i]),
\]
where $|V(\xi_i)|$ accounts for the choice of root within district $\xi_i$,
and $t(G[\xi_i])$ accounts for the choice of spanning tree within the district.

\paragraph{The stabiliser.}
Within $\mathrm{Pic}_k(G)$, the elements that map a forest compatible with
$\xi$ to another forest compatible with the same $\xi$ form the
\emph{stabiliser} $\mathrm{Stab}(\xi) \subset \mathrm{Pic}_k(G)$, with
\[
  |\mathrm{Stab}(\xi)| = \prod_{i=1}^{k} |V(\xi_i)| \cdot t(G[\xi_i]).
\]
The number of distinct partitions reachable from $\xi$ by the group action is:
\[
  |\mathrm{Pic}_k(G)| \,/\, |\mathrm{Stab}(\xi)|.
\]
If the action of $\mathrm{Pic}_k(G)$ on rooted $k$-forests is transitive,
this quotient equals the total number of connected $k$-partitions (with the
given root set).

\warning{The stabiliser is only defined if $\mathrm{Pic}_k(G)$ has a proven
group action on rooted $k$-forests. Currently this action is conjectured (based
on the counting match $|\mathrm{Pic}_k(G)| = $ number of rooted $k$-forests)
but not proven. Without the action, $\mathrm{Stab}(\xi)$ is undefined and the
formula $|\mathrm{Stab}(\xi)| = \prod |V(\xi_i)| \cdot t(G[\xi_i])$ is a
heuristic based on fibre counting, not a theorem. Proving this action (e.g.,
via a multi-sink Bernardi bijection or multi-sink rotor-routing torsor) is the
central open problem.}

\paragraph{Embedding into $\Pic(G)$.}
A rooted $k$-forest with roots $\{q_1, \ldots, q_k\}$ (where $q_1 = q$ is
the original sink of $\Pic(G)$) corresponds to a divisor in $\Pic(G)$ where
the coordinates at $q_2, \ldots, q_k$ are zero. Define the embedding
$\iota : \mathrm{Pic}_k(G) \hookrightarrow \Pic(G)$ by:
\[
  \iota(D)_v =
  \begin{cases}
    D_v & \text{if } v \notin \{q_2, \ldots, q_k\}, \\
    0   & \text{if } v \in \{q_2, \ldots, q_k\}.
  \end{cases}
\]
The image of $\iota$ is the subgroup of $\Pic(G)$ where the coordinates at
$q_2, \ldots, q_k$ vanish. The surjection
$\pi : \Pic(G) \to \mathrm{Pic}_k(G)$ (forgetting these coordinates) is the
left inverse: $\pi \circ \iota = \mathrm{id}$.

\warning{The map $\iota$ is NOT a well-defined group homomorphism. Equivalence
in $\mathrm{Pic}_k(G)$ allows chips to vanish at the $k-1$ extra sinks
$q_2, \ldots, q_k$. In $\Pic(G)$, those vertices are present and chips do not
vanish. So $D \sim D'$ in $\mathrm{Pic}_k(G)$ does NOT imply
$\iota(D) \sim \iota(D')$ in $\Pic(G)$. The map $\iota$ acts on
representatives, not on equivalence classes. The surjection
$\pi : \Pic(G) \to \mathrm{Pic}_k(G)$ IS well-defined (projecting respects
equivalence because $\operatorname{im}(\Delta^{(q_1)})$ projects onto a
superset of $\operatorname{im}(\Delta^{(\{q_1,\ldots,q_k\})})$), but the
reverse embedding does not exist as a group homomorphism.}

\paragraph{The nested structure.}
This gives a clean hierarchy:
\[
  \{k\text{-partitions}\}
    \;\subset\; \{k\text{-forests}\}
    \;\subset\; \mathrm{Pic}_k(G)
    \;\hookrightarrow\; \Pic(G)
    \;\longleftrightarrow\; \{\text{spanning trees}\}.
\]

\warning{This diagram mixes rigorous and conjectured relationships. The
$\longleftrightarrow$ (Baker--Wang torsor) is proven. The $\hookrightarrow$ is
NOT a valid group homomorphism (see warning above). The second $\subset$
($k$-forests $\subset \mathrm{Pic}_k(G)$) assumes a bijection between
$\mathrm{Pic}_k(G)$ and rooted $k$-forests, which is a counting match
($|\mathrm{Pic}_k(G)| = $ number of rooted $k$-forests via Matrix-Forest
Theorem) but NOT a proven group action or torsor. The first $\subset$
($k$-partitions $\subset k$-forests) is a set-level quotient, not an algebraic
inclusion.}

At each level:
\begin{itemize}
  \item $\Pic(G)$ has $t(G)$ elements, indexing spanning trees.
  \item $\mathrm{Pic}_k(G)$ has $\det(\Delta^{(\{q_1,\ldots,q_k\})})$
    elements, indexing rooted $k$-forests.
  \item $k$-forests modulo within-district equivalence give $k$-partitions.
\end{itemize}
Each inclusion has computable index. The group structure of $\Pic(G)$ governs
the entire hierarchy, and the partition structure is identified by the
stabiliser at each level.

\paragraph{Remark.}
Note that $\mathrm{Stab}(\xi)$ depends on $\xi$, so different partitions have
different stabilisers. However, since $\mathrm{Pic}_k(G)$ is abelian, every
subgroup is normal, and the quotient $\mathrm{Pic}_k(G) / \mathrm{Stab}(\xi)$
is a well-defined group for each $\xi$. The key open question is whether this
quotient is ``the same'' (isomorphic) for all $\xi$, which would give a
universal ``partition Picard group'' independent of the specific partition.
\end{answer}


%═══════════════════════════════════════════════════════════════════════════════
\section{Literature Review: Key Papers}
%═══════════════════════════════════════════════════════════════════════════════

\subsection{Torsor Structures on Spanning Trees}

\textbf{Baker and Wang}~\cite{baker2018bernardi} prove that for planar graphs,
the Bernardi process defines a \emph{simply transitive} action of $\Pic(G)$ on
spanning trees (Theorem~4.1). This action is canonical for planar graphs
(independent of root vertex, Theorem~5.1) and coincides with the rotor-routing
torsor (Theorem~7.1). This is the foundational result for the algebraic
approach to spanning trees.

\textbf{Backman, Baker, and Yuen}~\cite{backman2019geometric} extend this to
\emph{regular matroids}: the Jacobian of a regular matroid admits a canonical
simply transitive action on circuit-cocircuit reversal classes. The bijections
$\beta$ are constructed via acyclic signatures.

\textbf{Ding, McDonough, T\'othm\'er\'esz, and Yuen}~\cite{ding2025consistent}
prove that the sandpile torsor on regular matroids is \emph{consistent} under
deletion and contraction. This is the key property for our setting: contracting
within-district edges gives the quotient graph $Q_\xi$, and consistency means
the torsor on $G$ projects to a torsor on $Q_\xi$.

\subsection{Bijections for Spanning Subgraphs}

\textbf{Ding}~\cite{ding2023geometric} extends the BBY bijection from spanning
trees to \emph{all spanning subgraphs}. Theorem~1.2 gives specialisations:
\begin{enumerate}
  \item Spanning subgraphs $\leftrightarrow$ all orientations (cardinality
    $T_G(2,2)$).
  \item \textbf{Spanning forests} $\leftrightarrow$ $\sigma$-compatible
    orientations (cardinality $T_G(2,1)$).
  \item Connected spanning subgraphs $\leftrightarrow$ $\sigma^*$-compatible
    orientations (cardinality $T_G(1,2)$).
  \item Spanning trees $\leftrightarrow$ $(\sigma, \sigma^*)$-compatible
    orientations (cardinality $T_G(1,1)$).
\end{enumerate}
Row~2 is critical for us: spanning forests (= partitions into connected
subgraphs) biject with $\sigma$-compatible orientations. However, the group
action (torsor) is only established at the spanning tree level (row~4), not at
the forest level.

\subsection{Partitions and Divisors}

\textbf{Xu}~\cite{xu2011partitions} establishes a direct map from ordered
$k$-partitions to divisors. For an ordered partition $S = (V_1, \ldots, V_k)$,
the \emph{minimal effective divisor} $D^S$ has value at each vertex $v \in V_i$
equal to the number of neighbours of $v$ in districts $V_j$ with $j > i$
(Proposition~2). Key consequences:
\begin{itemize}
  \item Chip-firing $\chi(V_1)$ rotates the ordered partition:
    $\chi(V_1) \cdot D^S = D^{S^{(1)}}$ (Corollary~4).
  \item Merging adjacent districts $V_i \cup V_{i+1}$ corresponds to
    subtracting the boundary divisor: $D^S - D^{S'} = \sum_{v \in V_i}
    |\{vw \in E : w \in V_{i+1}\}| \cdot v$ (Corollary~3).
  \item The minimal free resolution of the toppling ideal $R/I$ is
    (conjecturally) indexed by connected partitions (Conjecture~13, due to
    Wilmes).
\end{itemize}

\subsection{The Matrix-Forest Theorem and $W(G) = I + L(G)$}

\textbf{Chebotarev and Shamis}~\cite{chebotarev2006forest} prove the
Matrix-Forest Theorems for the matrix $W(G) = I + L(G)$:
\begin{itemize}
  \item $\det(W(G))$ = total number of rooted spanning forests (Theorem~5).
  \item The $(i,j)$-cofactor $W^{ij}$ = number of rooted forests where $i$ and
    $j$ belong to the same tree rooted at $i$ (Theorem~6).
  \item $W^{-1}_{ij}$ = relative forest-accessibility: the fraction of forests
    where $i$ and $j$ are in the same tree rooted at $j$ (Theorem~7).
\end{itemize}
The inverse $W^{-1}$ is a \emph{partition proximity matrix}: vertices with
large $W^{-1}_{ij}$ naturally belong in the same district across random
forests. However, $W(G)$ counts forests of \emph{all} component numbers
simultaneously and does not directly encode $k$-partitions.

The referenced papers provide additional structure:
\begin{itemize}
  \item \textbf{Bapat and Constantine}~\cite{bapat1992enumerating} count
    spanning forests with colour restrictions --- potentially useful for
    encoding within-district vs.\ cross-boundary edge constraints algebraically.
  \item \textbf{Kelmans and Chelnokov}~\cite{kelmans1974polynomial} show that
    the coefficients of the Laplacian characteristic polynomial equal the
    numbers of spanning forests with fixed numbers of components. This connects
    the \emph{spectrum} of $\Delta_G$ to forest/partition counts at every level
    $k$ simultaneously.
  \item \textbf{Merris}~\cite{merris1989edge} proves an edge version of the
    Matrix-Tree Theorem using the edge Laplacian $B^T B$, potentially enabling
    an edge-based analogue of $\mathrm{Pic}_k(G)$.
  \item \textbf{Moon}~\cite{moon1995adjoint} studies the adjoint of $W(G)$,
    which by Theorem~7 of~\cite{chebotarev2006forest} gives the matrix of
    relative forest-accessibilities.
\end{itemize}

\subsection{Redistricting MCMC}

\textbf{DeFord, Herschlag, and Mattingly}~\cite{deford2025cycle} introduce
the Cycle Walk for sampling spanning forests for redistricting --- the closest
existing work to our approach.

\textbf{Autry, Carter, Herschlag, Hunter, and
Mattingly}~\cite{autry2020multiscale} develop multi-scale merge-split MCMC,
addressing the same scaling problem as our boundary cycle step but via
hierarchical proposals rather than algebraic structure.

\textbf{McWhorter and DeFord}~\cite{mew} introduce the Marked Edge Walk, whose
within-fibre waste motivates our entire programme.


\begin{thebibliography}{99}

\bibitem{autry2020multiscale}
E.~A. Autry, D.~Carter, G.~Herschlag, Z.~Hunter, and J.~C. Mattingly.
\newblock Multi-scale merge-split {M}arkov chain {M}onte {C}arlo for
  redistricting.
\newblock \emph{arXiv preprint arXiv:2008.08054}, 2020.

\bibitem{backman2019geometric}
S.~Backman, M.~Baker, and C.~H. Yuen.
\newblock Geometric bijections for regular matroids, zonotopes, and {E}hrhart
  theory.
\newblock \emph{Forum of Mathematics, Sigma}, 7:e45, 2019.

\bibitem{baker2018bernardi}
M.~Baker and Y.~Wang.
\newblock The {B}ernardi process and torsor structures on spanning trees.
\newblock \emph{International Mathematics Research Notices}, 2018(16):5120--5147,
  2018.

\bibitem{bapat1992enumerating}
R.~B. Bapat and G.~Constantine.
\newblock An enumerating function for spanning forests with color restrictions.
\newblock \emph{Linear Algebra Appl.}, 173:231--237, 1992.

\bibitem{chebotarev2006forest}
P.~Chebotarev and E.~Shamis.
\newblock Matrix-forest theorems.
\newblock \emph{arXiv preprint arXiv:math/0602575}, 2006.

\bibitem{deford2025cycle}
D.~R. DeFord, G.~Herschlag, and J.~C. Mattingly.
\newblock A cycle walk for sampling measures on spanning forests for
  redistricting.
\newblock \emph{arXiv preprint arXiv:2509.08629}, 2025.

\bibitem{dhar1990self}
D.~Dhar.
\newblock Self-organized critical state of sandpile automaton models.
\newblock \emph{Physical Review Letters}, 64(14):1613, 1990.

\bibitem{ding2023geometric}
C.~Ding.
\newblock Geometric bijections between spanning subgraphs and orientations of a
  graph.
\newblock \emph{arXiv preprint arXiv:2109.01930}, 2023.

\bibitem{ding2025consistent}
C.~Ding, A.~McDonough, L.~T\'othm\'er\'esz, and C.~H. Yuen.
\newblock A consistent sandpile torsor algorithm for regular matroids.
\newblock \emph{European Journal of Combinatorics}, 130:104218, 2025.

\bibitem{kelmans1974polynomial}
A.~K. Kelmans and V.~M. Chelnokov.
\newblock A certain polynomial of a graph and graphs with an extremal number of
  trees.
\newblock \emph{J.\ Comb.\ Theory, Ser.\ B}, 16:197--214, 1974.

\bibitem{merris1989edge}
R.~Merris.
\newblock An edge version of the {M}atrix-{T}ree {T}heorem and the {W}iener
  index.
\newblock \emph{Lin.\ Multilin.\ Alg.}, 25:291--296, 1989.

\bibitem{mew}
A.~McWhorter and D.~DeFord.
\newblock The {M}arked {E}dge {W}alk: A novel {MCMC} algorithm for sampling of
  graph partitions.
\newblock \emph{arXiv preprint arXiv:2510.17714}, 2025.

\bibitem{moon1995adjoint}
J.~W. Moon.
\newblock On the adjoint of a matrix associated with trees.
\newblock \emph{Lin.\ Multilin.\ Alg.}, 39:191--194, 1995.

\bibitem{xu2011partitions}
T.~Xu.
\newblock \emph{Graph Partitions and Free Resolutions of Toppling Ideals}.
\newblock Bachelor's thesis, Reed College, 2011.

\end{thebibliography}


\end{document}
